\title{DeepSeekMath: Pushing the Limits of Mathematical Reasoning in Open Language Models}

\begin{abstract}
The intricacies and structured demands of mathematical reasoning render it particularly challenging for language models. In this work, we present DeepSeekMath~7B, developed by further pre-training DeepSeek-Coder-Base-v1.5 7B on a dataset comprising 120B math-related tokens curated from Common Crawl, alongside natural language and programming data. Remarkably, DeepSeekMath~7B attains a 51.7\% score on the rigorous, competition-grade MATH benchmark, achieving a performance level comparable to models such as Gemini-Ultra and GPT-4, all without the aid of external toolkits or ensembling strategies. Additionally, the model reaches 60.9\% on the MATH benchmark using self-consistency across 64 generated samples. The model’s strength in mathematical reasoning arises from two principal innovations: firstly, the effective utilization of extensive open-access web resources via a highly refined data selection process, and secondly, the development of Group Relative Policy Optimization (GRPO)—an adapted form of Proximal Policy Optimization (PPO)—which not only bolsters mathematical reasoning but also improves PPO’s memory efficiency.
\end{abstract}

\section{Introduction}

Large language models (LLMs) have dramatically transformed methodologies for mathematical reasoning within artificial intelligence, driving notable progress on quantitative reasoning tasks \citep{MATH} as well as on geometry-specific benchmarks \citep{trinh2024solving}. In addition, these systems have demonstrated substantial utility in aiding humans as they tackle challenging mathematical tasks \citep{tao}. Nevertheless, the most advanced systems—including GPT-4 \citep{gpt4} and Gemini-Ultra \citep{gemini}—remain inaccessible to the public, with existing open-source alternatives exhibiting much lower levels of performance by comparison.

This paper presents DeepSeekMath, a language model tailored specifically for mathematical tasks, which notably surpasses the mathematical reasoning capabilities of other open-source systems and closely rivals GPT-4's results on established academic evaluations. To build DeepSeekMath, we assemble the DeepSeekMath Corpus—a vast, meticulously curated pre-training resource containing 120B mathematics-related tokens. These data are gathered from the Common Crawl (CC) through the use of a classifier leveraging the fastText framework \citep{joulin2016fasttext}. Initially, the classifier is trained on data from OpenWebMath \citep{openwebmath} as positive samples, while negative samples are drawn from a wide array of unrelated web content. With this setup, we mine additional mathematical data from CC, subjecting selected items to human verification and labeling. The refined dataset is then used to further train and enhance the classifier's precision. Our analysis demonstrates that the resultant large-scale corpus is of high caliber: the base model DeepSeekMath-Base 7B records a 64.2\% score on GSM8K \citep{gsm8k} and reaches 36.2\% on the MATH dataset \citep{MATH}, thereby surpassing the performance reported by Minerva 540B \citep{minerva}. Moreover, because the DeepSeekMath Corpus encompasses multiple languages, we observe gains in Chinese-language mathematical evaluation datasets \citep{wei2023cmath,agieval}. We propose that the techniques developed for mathematical corpus curation here can serve as a foundation for future research, highlighting considerable potential for subsequent advances in this domain.

DeepSeekMath-Base utilizes DeepSeek-Coder-Base-v1.5 7B \citep{deepseek-coder} as its foundational model, based on our observation that initializing with a model trained on code yields superior outcomes relative to commencing from a standard LLM. Additionally, our findings suggest that mathematical training enhances performance not only on mathematics-specific tasks, but also boosts results on the MMLU \citep{mmlu} and BBH \citep{bbh} benchmarks. This suggests that mathematical pretraining strengthens both mathematical proficiency and general reasoning skills.

Following the pre-training phase, we fine-tune DeepSeekMath-Base using mathematical instruction data that includes chain-of-thought \citep{cot}, program-of-thought \citep{pot,pal}, and tool-integrated reasoning \citep{tora} approaches. The fine-tuned model, DeepSeekMath-Instruct 7B, outperforms all other 7B-scale models and demonstrates performance on par with open-source instruction-tuned models at the 70B scale.

In addition, we propose Group Relative Policy Optimization (GRPO), a reinforcement learning (RL) algorithm that modifies Proximal Policy Optimization (PPO) \citep{schulman2017proximal}. Unlike conventional PPO, GRPO dispenses with the critic model by deriving the baseline directly from group scores, leading to a notable decrease in the computational cost during training. When utilizing only a fraction of English instruction tuning data, GRPO delivers notable gains over the robust DeepSeekMath-Instruct model, enhancing performance on both in-domain tasks (GSM8K: 82.9\% $\rightarrow$ 88.2\%, MATH: 46.8\% $\rightarrow$ 51.7\%) and on out-of-domain problems (for instance, CMATH: 84.6\% $\rightarrow$ 88.8\%) in the reinforcement learning stage. Additionally, we propose a unified framework for interpreting a variety of approaches, including Rejection Sampling Fine-Tuning (RFT) \citep{yuan2023scaling}, Direct Preference Optimization (DPO) \citep{dpo}, PPO, and GRPO. Through this unified perspective, we recognize that these algorithms may be interpreted as variations or simplifications of RL methods. We further carry out comprehensive analyses, such as comparing online and offline training regimes, assessing supervision at the outcome versus process level, and contrasting single-turn and iterative RL, in order to systematically dissect the key elements underlying this framework. Finally, we discuss the mechanisms by which our RL strategy enhances the effectiveness of instruction-tuned models, and we outline promising avenues for advancing RL methodology in line with the unified paradigm.

\subsection{Contributions}
We present scalable methods for math pre-training and conduct an in-depth investigation and evaluation of reinforcement learning.

\textbf{Math Pre-Training at Scale}

\begin{itemize}[topsep=0pt]
    \item This study demonstrates that the openly available Common Crawl dataset holds significant resources useful for mathematics. Through a carefully engineered pipeline for data selection, we assemble the DeepSeekMath Corpus—comprising 120B tokens sourced from web pages specifically screened for mathematical relevance. This collection is nearly sevenfold larger than the math-related web page data utilized by Minerva \citep{minerva}, and about nine times greater than the newly introduced OpenWebMath \citep{openwebmath}.

\item DeepSeekMath-Base 7B, our pre-trained foundational model, demonstrates performance on par with Minerva 540B \citep{minerva}. This suggests that parameter count alone does not determine proficiency in mathematical reasoning; a more compact model, when pre-trained using high-caliber datasets, can also exhibit robust capabilities.

\item We present the outcomes of our experiments involving mathematical training. Pretraining models on code before exposing them to mathematical tasks enhances their performance in solving math problems, regardless of whether tool use is involved. These results contribute to addressing the enduring inquiry: \textit{does code training enhance reasoning skills?} Our evidence indicates that it does, particularly in the realm of mathematical reasoning.

\item While utilizing arXiv papers as training data is a standard practice, particularly prevalent in numerous mathematics-focused studies, our findings reveal that this approach does not yield significant gains across any of the mathematical benchmarks employed in this work.
\end{itemize}

\textbf{Exploration and Analysis of Reinforcement Learning}

\begin{itemize}[topsep=0pt]
    \item We propose Group Relative Policy Optimization (GRPO), a novel and efficient reinforcement learning algorithm. Distinct from traditional approaches such as Proximal Policy Optimization (PPO), GRPO eliminates the need for a critic model by utilizing group score-based baseline estimation, resulting in substantial savings in computational resources required for training.
    \item Through empirical evaluation, we show that applying GRPO yields substantial improvements in the instruction-tuned model DeepSeekMath-Instruct, leveraging only instruction-tuning datasets. Moreover, reinforcement learning with GRPO contributes to better out-of-domain generalization.
    \item Our work offers a unified framework to interpret various existing methods—including RFT, DPO, PPO, and GRPO. Comprehensive experiments are performed to probe this framework, covering aspects like online versus offline training, outcome versus process-based supervision, and single-step compared to iterative reinforcement learning, among others, thereby elucidating the core components of the proposed paradigm.
    \item Leveraging this unified framework, we analyze underlying factors that drive the effectiveness of reinforcement learning, and highlight several promising avenues to further improve reinforcement learning approaches for large language models (LLMs).
\end{itemize}

\subsection{Summary of Evaluations and Metrics}

\begin{itemize}[topsep=0pt]
    \item \textbf{English and Chinese Mathematical Reasoning}: Our study entails a thorough evaluation of model performance on both English and Chinese mathematical reasoning benchmarks, which encompass a range of problems from primary school through university level. For English-language tasks, we utilize datasets such as GSM8K \citep{gsm8k}, MATH \citep{MATH}, SAT \citep{llemma}, OCW Courses \citep{minerva}, and MMLU-STEM \citep{mmlu}. For Chinese-language tasks, we include MGSM-zh \citep{mgsm}, CMATH \citep{wei2023cmath}, Gaokao-MathCloze \citep{agieval}, and Gaokao-MathQA \citep{agieval}. Model evaluation focuses on two key aspects: their capability to produce independent textual solutions unaided by external tools, and their proficiency in leveraging Python for mathematical problem solving.

On English evaluation datasets, DeepSeekMath-Base achieves results on par with the proprietary Minerva 540B model \citep{minerva}, and outperforms all publicly available base models, such as Mistral 7B \citep{mistral} and Llemma-34B \citep{llemma}, irrespective of their involvement in math-specific pre-training—often with a considerable lead. Remarkably, DeepSeekMath-Base attains even better outcomes on Chinese evaluation sets. This is plausibly due to our strategy of including not only English but also high-quality non-English mathematical data during pre-training, distinguishing our approach from that of previous studies \citep{minerva,llemma} that relied exclusively on English resources. Through mathematical instruction tuning and the use of reinforcement learning, the resultant DeepSeekMath-Instruct and DeepSeekMath-RL achieve impressive results, surpassing 50\% accuracy on the competition-grade MATH dataset—a first for open-source models.

\item \textbf{Formal Mathematics}: We assess DeepSeekMath-Base on the informal-to-formal theorem proving task from \citep{dsp_proof}, employing miniF2F \citep{minif2f} and selecting Isabelle \citep{isabelle} as the proof assistant. In this setting, DeepSeekMath-Base achieves robust performance in few-shot autoformalization scenarios.
\item \textbf{Natural Language Understanding, Reasoning, and Code}: To obtain a well-rounded evaluation of the models’ proficiency in general comprehension, logical reasoning, and coding, we conduct experiments with DeepSeekMath-Base on the Massive Multitask Language Understanding (MMLU) benchmark \citep{mmlu}, which features 57 multiple-choice tasks across a broad spectrum of topics; on BIG-Bench Hard (BBH) \citep{bbh}, a suite of 23 especially difficult tasks predominantly demanding multi-step reasoning; and on the established coding benchmarks HumanEval \citep{codex} and MBPP \citep{mbpp}. The mathematics-focused pre-training process positively influences both linguistic and reasoning competencies.

\section{Math Pre-Training}

\subsection{Data Collection and Decontamination} This section describes the methodology employed to assemble the DeepSeekMath Corpus utilizing data from Common Crawl. Figure \ref{fig:data_collect} illustrates an iterative workflow designed for systematically extracting an extensive mathematical dataset from Common Crawl, beginning with an initial high-quality seed corpus comprised of math-centric data. It is important to highlight that this methodology can be readily extended to additional fields, including but not limited to code-related corpora.

To begin, OpenWebMath \citep{openwebmath}—an extensive repository of mathematically-oriented, high-quality online texts—is utilized as the foundational seed corpus. We leverage this dataset to train a fastText model \citep{joulin2016fasttext} aimed at retrieving additional web pages with characteristics similar to those found in OpenWebMath. Specifically, 500,000 samples are randomly drawn from the seed corpus and labeled as positive examples, while an equivalent number of randomly sampled web pages from Common Crawl are designated as negative examples. Model training is performed using an open-source implementation\footnote{\url{https://fasttext.cc}}, setting the vector dimensionality at 256, specifying a learning rate of 0.1, a maximum word n-gram length of 3, a minimum word occurrence threshold of 3, and training over 3 epochs. To effectively narrow down the initial Common Crawl dataset, both URL-based deduplication and near-duplicate removal techniques are applied, reducing the dataset to 40 billion HTML web pages. The trained fastText model is then employed to identify mathematical web content from the deduplicated Common Crawl. Post-retrieval, the identified pages are ranked based on their model-predicted scores, with only the highest-scoring documents being retained to ensure quality. The resulting data volume is further evaluated by conducting pre-training runs on selections comprising the top 40B, 80B, 120B, and 160B tokens. For the initial phase, we opt to retain just the top 40B tokens.

Following the initial round of data harvesting, a considerable number of mathematical web pages are still missing from the collection. This shortfall can be attributed to the fastText model, which was initially trained on a collection of positive examples lacking in variety. To remedy this, we locate supplementary mathematical web domains to expand and diversify the seed corpus, thereby facilitating further optimization of the fastText model. Our methodology begins by dividing the complete set of Common Crawl data into separate domains, where each domain encompasses web pages with an identical base URL. For each of these domains, we compute the proportion of web pages captured in the first round of collection. Those domains where more than 10\% of web pages have been acquired are designated as being related to mathematics (for example, \url{mathoverflow.net}).

Next, we proceed with the manual annotation of URLs within these math-related domains that correspond to mathematical content (such as \url{mathoverflow.net/questions}). Any web pages associated with these URLs that have not yet been harvested are then included in the seed corpus. This methodology allows us to augment the pool of positive examples used for training, thus enhancing the fastText model’s ability to identify and retrieve a greater volume of mathematical content in subsequent rounds. After executing four such rounds, the corpus comprises 35.5M mathematical web pages, amounting to a total of 120B tokens. During the fourth round, we observe that about 98\% of the data had already been gathered by the previous round, prompting the decision to halt further data acquisition.

In order to prevent contamination of our benchmarks, we adopt the approach outlined in \cite{deepseek-coder}, filtering out web pages that include questions or answers sourced from well-known mathematical evaluation sets such as GSM8K \citep{gsm8k}, MATH \citep{MATH}, CMATH \citep{wei2023cmath}, and AGIEval \citep{agieval}. The filtering process is defined as follows: any segment of text that contains a 10-gram sequence that exactly matches a substring from the listed benchmarks is excluded from our math training data. For benchmark items with fewer than 10 grams but with a minimum of 3 grams, we apply precise matching to identify and remove potentially contaminated web content.

\subsection{Validating the Quality of the DeepSeekMath~Corpus}

We run pre-training experiments to investigate how the DeepSeekMath~Corpus is compared with the recently released math-training corpora:

\begin{itemize}[topsep=0pt]
    \item \textbf{MathPile} \citep{mathpile}: This dataset comprises 8.9B tokens collected from diverse sources such as textbooks, Wikipedia, ProofWiki, CommonCrawl, StackExchange, and arXiv. Notably, arXiv accounts for more than 85\% of the total tokens.
    \item \textbf{OpenWebMath} \citep{openwebmath}: This corpus includes 13.6B tokens and is constructed by extracting mathematical content from CommonCrawl.
    \item \textbf{Proof-Pile-2} \citep{llemma}: Assembled by merging OpenWebMath, AlgebraicStack (which contains 10.3B tokens of code relevant to mathematics), and papers from arXiv (28.0B tokens). For experiments involving Proof-Pile-2, we adhere to the 2:4:1 arXiv:Web:Code composition proposed in \cite{llemma}.
\end{itemize}

\subsubsection{Training Setting}
\label{sec:quality-policy}
Mathematics-specific training is performed on a broadly pre-trained language model comprising 1.3B parameters, constructed using the same architecture as the DeepSeek LLMs \citep{deepseek-llm}, referred to here as DeepSeek-LLM 1.3B. For each distinct mathematical dataset, an independent model undergoes training over 150B tokens. The HAI-LLM \citep{haillm} framework, noted for its efficiency and minimal computational overhead, underpins all of the experiments. Consistent with established DeepSeek LLMs training routines, the AdamW optimizer \citep{adamW} is utilized with $\beta_1=0.9$, $\beta_2=0.95$, and a $\mathrm{weight\_decay}$ of 0.1. The learning rate follows a multi-stage schedule: it ramps up to its maximum within the first 2,000 warmup steps, is reduced to 31.6\% of its maximum after completing 80\% of the total training, and is further curtailed to 10.0\% of the maximum at 90\% training progress. The learning rate’s maximum threshold is specified at $5.3 \times 10^{-4}$. Training batches encompass 4M tokens, and sequences are processed with a context length of 4K.

\subsubsection{Evaluation Results}

\textbf{The DeepSeekMath~Corpus is of high quality, covers multilingual mathematical content, and is the largest in size.}

\begin{itemize}[topsep=0pt]
    \item \textbf{High-quality}:
    We assess downstream performance on eight mathematical evaluation benchmarks using few-shot chain-of-thought prompting \cite{cot}. According to Table \ref{tab:corpora_comparison}, models trained with the DeepSeekMath~Corpus exhibit a marked advantage in performance. Moreover, Figure \ref{fig:corpora_comparisons} reveals that models trained on the DeepSeekMath Corpus outperform those trained on Proof-Pile-2 at the 50B token mark (which constitutes one complete epoch of Proof-Pile-2), demonstrating that the mean quality of DeepSeekMath Corpus surpasses that of existing alternatives.
    \item \textbf{Multilingual}:
    The DeepSeekMath~Corpus contains mathematical data in several languages, with English and Chinese making up the majority. As indicated in Table \ref{tab:corpora_comparison}, models trained on DeepSeekMath~Corpus achieve superior mathematical reasoning abilities in both English and Chinese contexts. In comparison, most other mathematical corpora are overwhelmingly English-focused, delivering limited gains and sometimes even negatively affecting Chinese mathematical reasoning performance.
    \item \textbf{Large-scale}:
    The DeepSeekMath~Corpus exceeds prior mathematical corpora by a substantial margin in terms of scale. Figure \ref{fig:corpora_comparisons} illustrates that DeepSeek-LLM 1.3B, when pre-trained with DeepSeekMath~Corpus, demonstrates a sharper improvement curve and sustained advancement. Conversely, baseline corpora are considerably smaller, undergo several repeat training cycles, and yield model performances that plateau rapidly.
\end{itemize}

\subsection{Training and Evaluating DeepSeekMath-Base 7B}

This section presents DeepSeekMath-Base 7B, a foundational model characterized by advanced reasoning skills, particularly in the area of mathematics. The model’s initial parameters are sourced from DeepSeek-Coder-Base-v1.5 7B \citep{deepseek-coder}, with subsequent training conducted over 500B tokens. The data utilized in training is composed as follows: DeepSeekMath Corpus comprises 56\%, AlgebraicStack contributes 4\%, arXiv accounts for 10\%, code from Github represents 20\%, and the final 10\% consists of natural language text in English and Chinese from Common Crawl. The majority of the training protocol adheres to the configuration described in Section \ref{sec:quality-policy}; however, the peak learning rate is set at 4.2e-4, and the batch size is established at 10M tokens.

We perform an extensive evaluation of DeepSeekMath-Base 7B's mathematical proficiency, examining its capacity to generate complete mathematical solutions independently, utilize tools for solving mathematical tasks, and engage in formal theorem proving. In addition to its mathematical functionalities, we present a broader overview of the base model's general abilities, assessing its competence in natural language understanding, logical reasoning, and programming.

\paragraph{Mathematical Problem Solving with Step-by-Step Reasoning}

We assess the ability of DeepSeekMath-Base to solve mathematical tasks through few-shot chain-of-thought prompting \citep{cot} on eight distinct English and Chinese benchmarks. These benchmarks include both quantitative reasoning tasks—such as GSM8K \citep{gsm8k}, MATH \citep{MATH}, and CMATH \citep{wei2023cmath}—and multiple-choice formats like MMLU-STEM \citep{mmlu} and Gaokao-MathQA \citep{agieval}. The selected datasets span a wide range of mathematical domains, incorporating content from elementary to collegiate levels.

As indicated in Table \ref{tab:base_math_cot}, DeepSeekMath-Base 7B demonstrates superior results across all eight evaluation benchmarks when compared to other open-source base models, which include the popular Mistral 7B \citep{mistral} and the newer Llemma 34B \citep{llemma}—the latter of which underwent mathematical training on Proof-Pile-2 \citep{llemma}. Particularly remarkable is DeepSeekMath-Base’s performance on the challenging MATH dataset, where it achieves over 10\% higher absolute scores relative to any other open-source base models. Moreover, it exceeds the performance of Minerva 540B \citep{minerva}, a substantially larger, closed-source base model based on PaLM \citep{palm} and specialized with further training on mathematical corpora, despite being only a fraction of its size.

\paragraph{Mathematical Problem Solving with Tool Use} Our assessment of mathematical reasoning enhanced with programmatic tools is conducted on the GSM8K and MATH benchmarks, utilizing few-shot program-of-thought prompting as outlined by \citet{pot,pal}. The approach requires models to address each question by generating Python code that may invoke libraries such as \textit{math} and \textit{sympy} for handling sophisticated mathematical procedures. The output produced by running the generated code is treated as the model’s solution. Table \ref{tab:base_math_pot_proof} demonstrates that DeepSeekMath-Base 7B achieves superior results compared to the earlier state-of-the-art model, Llemma 34B.

\paragraph{Formal Mathematics}

Automating formal proofs contributes significantly to the precision and dependability of mathematical argumentation, as well as operational efficiency, and has garnered notable interest in recent times. In this study, we assess DeepSeekMath-Base 7B on the informal-to-formal proof generation task from \citep{dsp_proof}, where the objective is to construct a formal proof using an informal description, the corresponding formal statement, and an informal argument. For this evaluation, we utilize the miniF2F \citep{minif2f} dataset, which serves as a benchmark for formalized mathematics at Olympiad level, prompting the model with a few examples to create formal proofs in Isabelle for each instance. In alignment with the methodology from \cite{dsp_proof}, proof outlines are generated by the model, and we then employ the existing automated prover Sledgehammer \citep{sledgehammer} to resolve any remaining proof steps. Results, as presented in Table \ref{tab:base_math_pot_proof}, indicate that DeepSeekMath-Base 7B performs robustly in the context of proof autoformalization.

\paragraph{Natural Language Understanding, Reasoning, and Code}

We assess the capabilities of models in natural language understanding using MMLU \citep{mmlu}, reasoning skills via BBH \citep{bbh}, and coding proficiency on HumanEval \citep{codex} and MBPP \citep{mbpp}. According to the results presented in Table \ref{tab:understanding_reasoning_code}, DeepSeekMath-Base 7B achieves notable improvements on MMLU and BBH compared to its predecessor, DeepSeek-Coder-Base-v1.5 \citep{deepseek-coder}, underscoring the beneficial effects of mathematical pretraining on both language comprehension and reasoning abilities. Furthermore, by continuing to expose the model to code tokens during ongoing training, DeepSeekMath-Base 7B successfully preserves the strong coding performance established by DeepSeek-Coder-Base-v1.5 on both coding-related benchmarks. In summary, DeepSeekMath-Base 7B yields superior results to the general-purpose Mistral 7B model \citep{mistral} across all three evaluated benchmarks in reasoning and coding.

\section{Supervised Fine-Tuning}

\subsection{SFT Data Curation}

We develop a dataset for mathematical instruction-tuning that encompasses both English and Chinese questions spanning a range of mathematical domains and difficulty levels. Each problem is matched with solutions represented in chain-of-thought (CoT) \citep{cot}, program-of-thought (PoT) \citep{pot,pal}, and tool-integrated reasoning formats \citep{tora}. The dataset comprises a total of 776,000 training examples.

\begin{itemize}[topsep=0pt]
    \item \textbf{English mathematical datasets}: Our work involves annotating GSM8K and MATH datasets with solutions that incorporate external tool usage, and additionally utilizing a subset of MathInstruct \citep{MathInstruct} as well as the Lila-OOD training set \citep{lila}, which feature problems addressed using CoT or PoT strategies. The English dataset encompasses a wide spectrum of mathematical domains, including but not limited to algebra, probability, number theory, calculus, and geometry.
    \item \textbf{Chinese mathematical datasets}: We compile K-12 level math questions in Chinese, which span 76 distinct subfields such as linear equations, and provide annotations in both CoT and tool-integrated reasoning formats.
\end{itemize}

\subsection{Training and Evaluating DeepSeekMath-Instruct 7B}

This section presents DeepSeekMath-Instruct 7B, a model that is refined through mathematical instruction tuning using DeepSeekMath-Base as its foundation. During training, examples are combined in random order until a total length of 4K tokens is achieved. The model is trained over 500 iterations with a batch size set at 256, employing a fixed learning rate of 5e-5.

We evaluate models' mathematical performance both without and with tool use, on 4 quantitative reasoning benchmarks in English and Chinese. We benchmark our model against the leading models of the time:

\begin{itemize}[topsep=0pt]
    \item \textbf{Closed-source models} comprise: (1) the GPT series, particularly GPT-4 \citep{gpt4} and GPT-4 Code Interpreter~\footnote{\url{https://openai.com/blog/chatgpt-plugins\#\#code-interpreter}} as leading representatives, (2) Gemini Ultra and Pro \citep{gemini}, (3) Inflection-2 \citep{inflection-2}, (4) Grok-1~\footnote{\url{https://x.ai/model-card}},
    and several models introduced by Chinese enterprises such as (5) Baichuan-3~\footnote{\url{https://www.baichuan-ai.com}},
    and (6) the most recent GLM-4~\footnote{\url{https://open.bigmodel.cn/dev/api\#glm-4}}

    from the GLM family \citep{glm}.

These models are intended for broad applications, and most have been subjected to various alignment strategies.

\item \textbf{Open-source models} encompass both general-purpose and mathematically enhanced variants: among the former are (1) DeepSeek-LLM-Chat 67B \citep{deepseek-llm}, (2) Qwen 72B \citep{qwen}, (3) SeaLLM-v2 7B \citep{seallm}, and (4) ChatGLM3 6B \citep{chatglm3}. For mathematics-focused models, (5) InternLM2-Math 20B~\footnote{\url{https://github.com/InternLM/InternLM-Math}} extends InternLM2, receiving specialized math pretraining and subsequent instruction fine-tuning, (6) Math-Shepherd-Mistral 7B applies PPO optimization \citep{schulman2017proximal} to Mistral 7B \citep{mistral} guided by a process-supervised reward system, (7) the WizardMath collection \citep{wizardmath} enhances the math capabilities of Mistral 7B and Llama-2 70B \citep{llama2} via evolve-instruct—a technique for AI-generated instruction tuning—and PPO, primarily leveraging training datasets from GSM8K and MATH, (8) MetaMath 70B \citep{metamath}, where Llama-2 70B is further trained on enriched GSM8K and MATH corpora, (9) ToRA 34B \cite{tora}, representing CodeLlama 34B adapted for mathematical reasoning through tool augmentation, and (10) MAmmoTH 70B \citep{MathInstruct}, which comprises Llama-2 70B refined using MathInstruct-based instruction tuning.
\end{itemize}

As indicated in Table \ref{tab:sft_rl_math}, when evaluated without allowing tool use, DeepSeekMath-Instruct 7B exhibits impressive capabilities in step-by-step mathematical reasoning. In particular, on the competition-level MATH dataset, our model outperforms all available open-source systems as well as most proprietary alternatives—including Inflection-2 and Gemini Pro—by a margin of at least 9\% absolute. This superior performance persists even when compared with much larger models, such as Qwen 72B, and models that have received dedicated reinforcement learning improvements targeting mathematical problem solving, like WizardMath-v1.1 7B. Although DeepSeekMath-Instruct's results on MATH are on par with the Chinese proprietary models GLM-4 and Baichuan-3, the model does not yet reach the proficiency of leading models such as GPT-4 or Gemini Ultra.

When evaluated in a setup that permits models to combine natural language reasoning with programmatic tool utilization for solving problems, DeepSeekMath-Instruct 7B nearly attains a 60\% accuracy rate on the MATH benchmark, outperforming every previously available open-source model. Across additional benchmarks, our model demonstrates performance on par with DeepSeek-LLM-Chat 67B, the former state-of-the-art despite being ten times its size.
\section{Reinforcement Learning}

\subsection{Group Relative Policy Optimization} Reinforcement learning (RL) has demonstrated its value in enhancing the mathematical reasoning capabilities of LLMs beyond what is achieved through the Supervised Fine-Tuning (SFT) process \citep{wang2023math,wizardmath}. Here, we present Group Relative Policy Optimization (GRPO), a novel RL algorithm designed to be both efficient and effective.

\subsubsection{From PPO to GRPO} Proximal Policy Optimization (PPO) \citep{schulman2017proximal} is an actor-critic RL algorithm that is widely used in the RL fine-tuning stage of LLMs \citep{ouyang2022training}. In particular, it optimizes LLMs by maximizing the following surrogate objective:

\begin{equation}
\footnotesize
    \mathcal{J}_{PPO}(\theta) = \mathbb{E}{[q \sim P(Q), o \sim \pi_{\theta_{old}}(O|q)]} \frac{1}{|o|} \sum_{t=1}^{|o|} \min \left[ \frac{\pi_\theta(o_{t} | q, o_{<t})}{\pi_{\theta_{old}}(o_{t} | q, o_{<t})} A_{t}, \text{clip} \left( \frac{\pi_\theta(o_{t} | q, o_{<t})}{\pi_{\theta_{old}}(o_{t} | q, o_{<t})}, 1 - \varepsilon, 1 + \varepsilon \right)  A_{t} \right] ,
\end{equation}

The objective function for PPO, $\mathcal{J}_{PPO}(\theta)$, is defined as the expected value over prompts $q$ sampled from $P(Q)$ and output sequences $o$ drawn from the old policy $\pi_{\theta_{old}}(O|q)$. For each time step $t$ in the sequence $o$, the loss takes the average across all positions. At each position, it selects the smaller value between the importance sampling ratio times the advantage term, $A_t$, and a version of this ratio that has been clipped to the range $[1 - \varepsilon, 1 + \varepsilon]$, again multiplied by $A_t$. This approach ensures stable updates by penalizing large deviations from the previous policy.

Here, $\pi_{\theta}$ denotes the current policy, while $\pi_{\theta_{old}}$ represents the previous policy model. The variables $q$ and $o$ correspond to questions and their respective outputs, which are drawn from the question corpus and generated using the old policy $\pi_{\theta_{old}}$. The scalar $\varepsilon$ refers to a clipping hyper-parameter, implemented in PPO to promote stable training dynamics. The quantity $A_t$, representing the advantage, is estimated using Generalized Advantage Estimation (GAE) \citep{gae}; it relies on the sequence of rewards $\{r_{\ge t}\}$ and a learned value function $V_{\psi}$. Accordingly, PPO necessitates simultaneous training of both the policy and value networks. To address potential over-optimization of the reward model, it is customary to incorporate a KL-divergence penalty between the policy and a reference model at the token level within the reward computation \citep{ouyang2022training}, such that

\begin{equation}
    r_{t} = r_\varphi(q, o_{\le t}) - \beta \log\frac{\pi_{\theta}(o_{t}|q, o_{<t})}{\pi_{ref}(o_{t}|q, o_{<t})},
\label{eq:PPO-reward}
\end{equation}

Here, $r_\varphi$ denotes the reward model, $\pi_{ref}$ represents the reference model—typically the original SFT model—and $\beta$ corresponds to the weight applied to the KL divergence penalty.

As the value function employed in PPO is typically another model of comparable size as the policy model, it brings a substantial memory and computational burden. Additionally, during RL training, the value function is treated as a baseline in the calculation of the advantage for variance reduction. While in the LLM context, usually only the last token is assigned a reward score by the reward model, which may complicate the training of a value function that is accurate at each token. To address this, as shown in Figure \ref{fig:grpo}, we propose Group Relative Policy Optimization (GRPO), which obviates the need for additional value function approximation as in PPO, and instead uses the average reward of multiple sampled outputs, produced in response to the same question, as the baseline. More specifically, for each question $q$, GRPO samples a group of outputs $\{o_1, o_2, \cdots, o_G\}$  from the old policy  $\pi_{\theta_{old}}$  and then optimizes the policy model by maximizing the following objective:

\begin{equation}
\footnotesize
\begin{split}
    \mathcal{J}_{GRPO}(\theta) &= \mathbb{E}{[q \sim P(Q), \{o_i\}_{i=1}^G \sim \pi_{\theta_{old}}(O|q)]}  \\
    & \frac{1}{G}\sum_{i=1}^G\frac{1}{|o_i|} \sum_{t=1}^{|o_i|} \bigg\{ \min \left[ \frac{\pi_\theta(o_{i,t} | q, o_{i,<t})}{\pi_{\theta_{old}}(o_{i,t} | q, o_{i,<t})} \hat{A}_{i,t}, \text{clip} \left( \frac{\pi_\theta(o_{i,t} | q, o_{i,<t})}{\pi_{\theta_{old}}(o_{i,t} | q, o_{i,<t})}, 1 - \varepsilon, 1 + \varepsilon \right)  \hat{A}_{i,t} \right] \\
    & \qquad - \beta \mathbb{D}_{KL}\left[\pi_{\theta} || \pi_{ref}\right] \bigg\} ,
\end{split}
\label{eq:GRPO-obj}
\end{equation}

where $\varepsilon$ and $\beta$ are hyper-parameters, and $\hat{A}_{i,t}$  is the advantage calculated based on relative rewards of the outputs inside each group only, which will be detailed in the following subsections. The group relative way that GRPO leverages to calculate the advantages, aligns well with the comparative nature of rewards models,  as reward models are typically trained on datasets of comparisons between outputs on the same question. Also note that, instead of adding KL penalty in the reward, GRPO regularizes by directly adding the KL divergence between the trained policy and the reference policy to the loss, avoiding complicating the calculation of  $\hat{A}_{i,t}$. And different from the KL penalty term used in (\ref{eq:PPO-reward}), we estimate the KL divergence with the following unbiased estimator \citep{kl_approx}:

\begin{equation}
\small
    \mathbb{D}_{KL}\left[\pi_{\theta} || \pi_{ref}\right] = \frac{\pi_{ref}(o_{i,t}|q,o_{i,<t})}{\pi_{\theta}(o_{i,t}|q,o_{i,<t})} - \log\left(\frac{\pi_{ref}(o_{i,t}|q,o_{i,<t})}{\pi_{\theta}(o_{i,t}|q,o_{i,<t})}\right) - 1,
\end{equation}

which is guaranteed to be positive.

\begin{algorithm}[t]
  \small
  \caption{Iterative Group Relative Policy Optimization}
  \textbf{Input} initial policy model $\pi_{\theta_{\text{init}}}$; reward models $r_\varphi$; task prompts $\mathcal{D}$; 
  hyperparameters $\varepsilon$, $\beta$, $\mu$
  \begin{algorithmic}[1]
    \State Initialize the policy model as $\pi_\theta \leftarrow \pi_{\theta_{\text{init}}}$
    \For{iteration = 1, \dots, I}
       \State Set the reference model to current policy: $\pi_{ref} \leftarrow \pi_{\theta}$
      \For{step = 1, \dots, M}
      \State Draw a mini-batch $\mathcal{D}_b$ from $\mathcal{D}$
      \State Save the present policy as $\pi_{\theta_{old}} \leftarrow \pi_{\theta}$ 
      \State For each query $q \in \mathcal{D}_b$, generate $G$ responses $\{o_i\}_{i=1}^G \sim \pi_{\theta_{old}} (\cdot \mid q)$
      \State For every response $o_i$, determine the rewards $\{r_i\}_{i=1}^{G}$ using the reward model $r_{\varphi}$ 
      \State Estimate the group relative advantage $\hat{A}_{i,t}$ for each token $t$ within $o_i$
      \For{GRPO iteration = 1, \dots, $\mu$}
        \State Optimize the policy $\pi_{\theta}$ by maximizing the GRPO loss (see Equation \ref{eq:GC-GRPO})
      \EndFor
    \EndFor 
    \State Update the parameters of $r_\varphi$ via ongoing training with a replay buffer. 
    \EndFor 
  \end{algorithmic}
  \textbf{Output} $\pi_\theta$
  \label{alg:iter-grpo}
\end{algorithm}

\subsubsection{Outcome Supervision RL with GRPO} To describe this process formally, consider a question $q$, for which a collection of outputs $\{o_1, o_2, \cdots, o_G\}$ is drawn from the prior policy $\pi_{\theta_{old}}$. Each output in this group is evaluated using a reward model, producing a set of rewards $\mathbf{r}=\{r_1, r_2, \cdots, r_G\}$. These reward values are standardized by removing the group's mean and dividing by its standard deviation. Using outcome supervision, only the terminal normalized reward is supplied for each output $o_i$, and the advantage estimates $\hat{A}_{i, t}$ for every token in that output are set equal to this normalized reward: $\hat{A}_{i, t} = \widetilde{r}_i = \frac{r_i- {\rm mean}(\mathbf{r})}{{\rm std}(\mathbf{r})}$. The policy is then updated by maximizing the objective specified in equation (\ref{eq:GRPO-obj}).

\subsubsection{Process Supervision RL with GRPO}

Outcome supervision confines feedback to a single reward delivered at the conclusion of each generated output, which often proves insufficient and inefficient when guiding the policy for intricate mathematical reasoning. Building on the approach in \cite{wang2023math}, we implement process supervision, supplying a reward signal after each intermediate reasoning step. In this setup, for a given question $q$ and a set of $G$ sampled output sequences $\{o_1, o_2, \cdots, o_G\}$, a dedicated process reward model assesses each intermediate step and assigns step-wise reward values. This produces a set of rewards $\mathbf{R} = \{ \{r_1^{index(1)},\cdots,r_1^{index(K_1)}\}, \cdots,  \{r_G^{index(1)},\cdots,r_G^{index(K_G)}\} \}$, where $index(j)$ identifies the final token of the $j$-th reasoning step and $K_i$ counts the number of steps in the $i$-th sampled sequence. To standardize these signals, rewards are normalized using their empirical mean and standard deviation: $\widetilde{r}_i^{index(j)} = \frac{r_i^{index(j)} - {\rm mean(\mathbf{R})}}{{\rm std(\mathbf{R})}}$.

In the subsequent stage, process supervision computes the advantage for each token position as the cumulative sum of the normalized future step rewards: $\hat{A}_{i, t} = \sum_{index(j) \ge t} \widetilde{r}_i^{index(j)}$. The training process then updates the policy parameters by maximizing the objective function stated in equation (\ref{eq:GRPO-obj}).

\subsubsection{Iterative RL with GRPO} During reinforcement learning, the initially trained reward model may become inadequate for effectively guiding the evolving policy model. To address this, we investigate an iterative approach to RL utilizing GRPO. As described in Algorithm \ref{alg:iter-grpo}, this iterative GRPO framework involves constructing fresh training datasets for the reward model using the latest outputs sampled from the policy model. The reward model is then updated in an ongoing manner via a replay strategy that retains 10\% of prior data. Subsequently, the policy model is updated iteratively by adopting the latest policy as the reference model and optimizing it with the updated reward model.

\subsection{Training and Evaluating DeepSeekMath-RL}

We implement reinforcement learning (RL) utilizing DeepSeekMath-Instruct 7B as the base model. The RL training dataset comprises approximately 144,000 chain-of-thought-style questions from GSM8K and MATH, drawn exclusively from SFT data. Other SFT questions are deliberately omitted to specifically analyze the effects of RL on benchmarks for which no data is used during the RL stage. The reward model’s training dataset is curated in accordance with \citep{wang2023math}. Our reward model initialization employs DeepSeekMath-Base 7B, set with a learning rate of 2e-5. For GRPO, the policy model learning rate is maintained at 1e-6, and the KL divergence coefficient is fixed at 0.04. Each input question is used to generate $64$ sampled responses. The maximum sequence length is established at 1024, with a batch size also of 1024. After each exploration round, the policy model undergoes a single update. The evaluation of DeepSeekMath-RL 7B on benchmarks mirrors the procedure for DeepSeekMath-Instruct 7B. In this setup, GSM8K and MATH with chain-of-thought responses are classified as in-domain tasks, whereas the remaining benchmarks are considered out-of-domain.

Table \ref{tab:sft_rl_math} presents a comparison of both open- and closed-source models employing chain-of-thought and tool-assisted reasoning on English and Chinese evaluation sets. Our observations reveal: 1) When applying chain-of-thought reasoning, DeepSeekMath-RL 7B achieves accuracy rates of 88.2\% on GSM8K and 51.7\% on MATH. These results not only exceed the scores obtained by all open-source models in the 7B to 70B parameter spectrum but also outperform most closed-source competitors. 2) Notably, DeepSeekMath-RL 7B is exclusively trained on instruction tuning data formatted in chain-of-thought from GSM8K and MATH, initialized from DeepSeekMath-Instruct 7B. Even with this limited training data, it consistently outperforms DeepSeekMath-Instruct 7B across every evaluation dimension, thereby highlighting the advantages conferred by reinforcement learning.

\section{Discussion}
Here, we present the insights gained from our experiments involving pre-training and reinforcement learning (RL).

\subsection{Lessons Learnt in Pre-Training}

We begin by detailing our observations during pre-training. Unless stated otherwise, the experimental setup will follow the configurations described in Section \ref{sec:quality-policy}. Importantly, in this section, any mention of the DeepSeekMath Corpus pertains to an 89B-token dataset derived from the second round of data gathering.

\subsubsection{Code Training Benefits Mathematical Reasoning}

An often-cited but largely unproven assumption posits that code training enhances reasoning abilities. In this work, we seek to address this claim in part, focusing on mathematics: code training leads to enhanced mathematical reasoning skills in models, irrespective of whether auxiliary tools are employed.

To study how code training affects mathematical reasoning, we experimented with the following two-stage training and one-stage training settings:

\textbf{Two-Stage Training}

\begin{itemize}[topsep=0pt]
    \item \textbf{Code Training for 400B Tokens $\rightarrow$ Math Training for 150B Tokens}: The DeepSeek-LLM 1.3B model undergoes an initial training phase on 400B code tokens, which is then succeeded by an additional 150B math tokens for further training;
    \item \textbf{General Training for 400B Tokens $\rightarrow$ Math Training for 150B Tokens}: As a baseline for comparison, we substitute the 400B code tokens with general-purpose tokens (drawn from an extensive general domain corpus curated by DeepSeek-AI) in the first training stage, aiming to evaluate whether training on code tokens specifically confers greater benefits for mathematical reasoning relative to more general data.
\end{itemize}

\textbf{One-Stage Training}

\begin{itemize}[topsep=0pt]
    \item \textbf{Mathematical Training Using 150B Tokens}: The DeepSeek-LLM 1.3B model undergoes training with 150B tokens sourced from mathematical datasets;
    \item \textbf{Combined Training Utilizing 400B Code Tokens and 150B Math Tokens}: Sequential math training after initial code training results in diminished coding abilities. To address this, we explore whether concurrently training with both code and math tokens in a single-stage process can enhance the model's mathematical reasoning, while simultaneously mitigating the effects of catastrophic forgetting.
\end{itemize}

\paragraph{Results}
The downstream outcomes corresponding to various training configurations are displayed in Table \ref{tab:code-math} and Table \ref{tab:code-math-general-eval}.

Training with code bolsters program-aided mathematical reasoning capabilities, irrespective of whether a two-stage or one-stage training scheme is employed. Table \ref{tab:code-math} demonstrates that, in the context of two-stage training, exposure to code alone already leads to marked improvements in solving GSM8K and MATH tasks using Python. Incorporating math-specific training in the subsequent stage brings about additional gains. Notably, within the one-stage training regime, blending code tokens with math tokens both counteracts the issue of catastrophic forgetting observed in the two-stage approach, and produces mutually beneficial effects for both code-related tasks (Table \ref{tab:code-math-general-eval}) and program-aided mathematical reasoning (Table \ref{tab:code-math}).

Engaging in code training yields benefits for mathematical reasoning even in the absence of tool use. Within the framework of two-stage training, the preliminary phase focused on code already provides noticeable gains. This phase further increases the effectiveness of the following math-oriented training, ultimately achieving superior results. In contrast, merging code and math tokens in a single-stage training paradigm diminishes mathematical reasoning performance when tools are not employed. A possible explanation is that DeepSeek-LLM 1.3B, constrained by its relatively small scale, may not possess sufficient capacity to effectively learn from both code and mathematical sources concurrently.

\subsubsection{ArXiv Papers Seem Ineffective in Improving Mathematical Reasoning}

ArXiv papers are commonly included as a component of math pre-training data \citep{minerva,gpt-f,llemma,mathpile}. However, detailed analysis regarding their impact on mathematical reasoning has not been extensively conducted. Perhaps counter-intuitively, according to our experiments, arXiv papers seem ineffective in improving mathematical reasoning. We experiment with models of different sizes, including DeepSeek-LLM 1.3B and DeepSeek-Coder-Base-v1.5 7B \citep{deepseek-coder}, using arXiv corpora that underwent varied processing pipelines:

\begin{itemize}[topsep=0pt]
    \item \textbf{MathPile} \citep{mathpile}: a dataset comprising 8.9 billion tokens, assembled using a set of cleaning and filtering heuristics, with scientific arXiv articles representing over 85\% of its contents;
    \item \textbf{ArXiv-RedPajama} \citep{redpajama}: includes all arXiv LaTeX documents after excluding preambles, comments, macros, and references, resulting in a collection of 28.0 billion tokens.
\end{itemize}

In the conducted experiments, we individually trained DeepSeek-LLM 1.3B for 150B tokens and DeepSeek-Coder-Base-v1.5 7B for 40B tokens on distinct arXiv datasets. The results indicate that arXiv papers do not substantially enhance mathematical reasoning capabilities. When the models are trained exclusively on arXiv content, there is either negligible improvement or an actual decline observed in their performance across several mathematical benchmarks of varying difficulty examined in this work. These benchmarks comprise quantitative reasoning tasks such as GSM8K and MATH (Table \ref{tab:arxiv-cot}), multiple-choice assessments like MMLU-STEM (Table \ref{tab:arxiv-cot}), as well as formal mathematics evaluation via miniF2F (Table \ref{tab:arxiv_atp}).

However, this conclusion has its limitations and should be taken with a grain of salt. We have not yet studied:

\begin{itemize}[topsep=0pt]
    \item The influence of arXiv tokens on additional mathematical applications beyond the scope of this work, for instance, the informalization of theorems, which involves translating formal expressions or proofs into their informal counterparts;
    \item The interaction between arXiv tokens and other datasets, and its resultant impact;
    \item The extent to which advantages derived from arXiv papers persist or evolve as models increase in scale.
\end{itemize}

Thus, further exploration is required, which we leave for future studies.

\subsection{Insights of Reinforcement Learning}
\subsubsection{Towards to a Unified Paradigm} Here, we introduce a unified paradigm to systematically examine various training approaches, including SFT, RFT, DPO, PPO, and GRPO. We also design experiments to investigate the contributing factors of this unified framework. Broadly, the gradient of a training method with respect to the parameter $\theta$ can be expressed as:

\begin{equation}
    \nabla_{\theta}\mathcal{J}_{\textcolor{red}{\mathcal{A}}}(\theta) = \mathbb{E}[\underbrace{(q,o) \sim \textcolor{red}{\mathcal{D}}}_{Data \ Source}]\left( \frac{1}{|o|} \sum_{t=1}^{|o|}  \underbrace{GC_{{\mathcal{A}}}(q, o, t, \textcolor{red}{\pi_{{rf}}})}_{Gradient \ Coefficient}  \nabla_{\theta}\log \pi_{\theta}(o_t | q, o_{<t})\right).
\label{eq:objective}
\end{equation}

There exist three key components: 1) \textit{Data Source $\mathcal{D}$}, which determines the training data; 2) \textit{Reward Function $\pi_{{rf}}$}, which is the source of the training reward signal; 3) \textit{Algorithm $\mathcal{A}$}: which processes the training data and the reward signal to the gradient coefficient $GC$ that determines the magnitude of the penalty or reinforcement for the data. We analyze several representative methods based on such a unified paradigm:

\begin{itemize}
    \item \textbf{Supervised Fine-tuning (SFT)}: In SFT, the pretrained model is further adjusted using a dataset that has been carefully curated and labeled by humans.
    \item \textbf{Rejection Sampling Fine-tuning (RFT)}: RFT builds on the SFT model by further fine-tuning with a subset of responses generated from the SFT model itself in response to SFT prompts; these responses are selected based on the accuracy of their answers.
    \item \textbf{Direct Preference Optimization (DPO)}: DPO further tunes the SFT model by leveraging extra outputs, produced by the SFT model, and applying a pairwise DPO loss to enhance performance.
    \item \textbf{Online Rejection Sampling Fine-tuning (Online RFT)}: Online RFT diverges from the standard RFT in that it starts with the SFT model as the policy initialization, and continually refines the model using augmented outputs that are generated by the evolving, online policy model.
    \item \textbf{PPO/GRPO}: For PPO/GRPO, the policy model is initialized using weights from the SFT model, and its behavior is shaped further by sampling and reinforcing on real-time outputs from the current policy.
\end{itemize}

Table \ref{tab:data_policy} provides an overview of the elements comprising these methods. For a comprehensive explanation of the derivation steps, see Appendix \ref{app:analysis-rl}.

\paragraph{Observation about Data Source} The categorization of data sources is split into two types: online sampling and offline sampling. Online sampling refers to obtaining training data generated by the ongoing exploration of the currently used training policy model. In contrast, offline sampling involves collecting training data produced by the initial SFT model. Methods such as RFT and DPO utilize offline sampling, whereas Online RFT and GRPO adopt an online sampling approach.

As depicted in Figure \ref{fig:rl-analysis}, our results indicate that Online RFT achieves notably higher performance than RFT across both benchmarks. In the early training phase, the performance of Online RFT closely matches that of RFT; however, as training progresses, Online RFT gains a substantial lead, highlighting the benefits of online training methods. This outcome aligns with intuition: during the initial stages, the actor and SFT model are quite similar, so the data sampled reveals only slight variations. As training advances, the discrepancies between data from the actor become more pronounced, and the advantages of real-time data sampling become increasingly apparent.

\paragraph{Observation about Gradient Coefficient} The algorithm utilizes the reward signal to determine the gradient coefficient, which then guides parameter updates in the model. In our experimental framework, we separate the reward function into two categories: `Rule' and `Model'. The `Rule' reward is derived from evaluating answer quality purely in terms of correctness, while the `Model' reward is obtained by training a dedicated reward model to assign scores to responses. Notably, this reward model is supervised using data labeled according to the rule-based judgments. As shown in Equations \ref{eq:GC-RFT} and \ref{eq:GC-GRPO}, there is a fundamental distinction between GRPO and Online RFT. Specifically, GRPO tailors its gradient coefficient directly in response to the output of the reward model, enabling it to apply variable degrees of encouragement or discouragement that match the reward's intensity. On the other hand, Online RFT does not implement such nuanced feedback: it fails to penalize incorrect outputs, and consistently applies the same degree of reinforcement for all responses that are deemed correct, regardless of their individual quality.

As illustrated in Figure \ref{fig:rl-analysis}, GRPO outperforms online RFT, which underscores the effectiveness of adjusting both positive and negative gradient coefficients. Moreover, GRPO+PS achieves better results than GRPO+OS, demonstrating the advantage of employing fine-grained, step-specific gradient coefficients. Additionally, we investigate iterative RL by conducting two iterations in our experiments. Figure \ref{fig:iter-rl} reveals that iterative RL substantially enhances performance, with the most pronounced gains observed in the initial iteration.

\subsubsection{Why RL Works?} In this study, we apply reinforcement learning to a selected portion of instruction tuning datasets and observe that it markedly improves the performance relative to models trained solely through instruction tuning. To provide further insight into the mechanism behind reinforcement learning's effectiveness, we compare the Pass@K and Maj@K accuracy metrics for both Instruct and RL models across two evaluation benchmarks. As illustrated in Figure \ref{fig:maj-pass}, reinforcement learning yields a notable increase in Maj@K accuracy but does not affect Pass@K. This pattern suggests that RL's contribution lies in producing a more stable output distribution, meaning \textbf{the observed gains can largely be traced to increasing the frequency with which correct answers appear among the TopK outputs rather than to substantial improvement in underlying model competencies.} This observation is echoed by findings in \citep{wang2023making}, which document a \textbf{misalignment problem} in SFT models on reasoning tasks and demonstrate that such models can benefit from preference alignment approaches \citep{yuan2023rrhf,song2023preference,wang2023making}.

\subsubsection{How to Achieve More Effective RL?}
\label{sec:effec_rl}
Our findings indicate that RL is highly effective for mathematical reasoning tasks. Additionally, we introduce a comprehensive framework that encapsulates various notable training approaches, characterizing them as either direct or approximated forms of RL. According to Equation \ref{eq:objective}, this unified framework is comprised of three fundamental elements: Data Source, Algorithm, and Reward Function. We outline possible avenues for future research concerning each of these components.

\paragraph{Data Source} The data source forms the foundational input for any training methodology. Within reinforcement learning (RL), we define the data source as consisting of unlabeled questions paired with outputs that have been sampled from the policy model. For our experiments, we limit ourselves to the questions employed during the instruction tuning phase, utilizing basic nucleus sampling for generating outputs. We hypothesize that this restriction is a key factor behind our RL pipeline's improvements being primarily confined to Maj@K metrics. Looking forward, we intend to assess the efficacy of our RL pipeline on question prompts that lie outside the original distribution, alongside adopting \textbf{advanced sampling (decoding) strategies}, including those leveraging tree-search approaches \citep{yao2023tree}. Furthermore, recent developments in \textbf{efficient inference techniques} \citep{xia-etal-2023-speculative,leviathan2023fast,kwon2023efficient,xia2024unlocking}—which directly influence the effectiveness with which policy models explore—are expected to have significant impact as well.

\paragraph{Algorithms} Algorithms utilize both data and reward signals to generate gradient coefficients, which in turn are used to update the model parameters. As described in Equation \ref{eq:objective}, contemporary approaches typically place full \textbf{TRUST} in the reward signal to adjust the conditional probability of particular tokens, either increasing or decreasing it. Nevertheless, reward signals are not always dependable, particularly in tasks with high complexity. To illustrate, the PRM800K datasets \citep{lightman2023let}, despite being meticulously labeled by expert annotators, still exhibit around 20\% annotation errors\footnote{\url{https://github.com/openai/prm800k/issues/12\#issuecomment-1728491852}}. Given this limitation, it becomes imperative to investigate reinforcement learning algorithms that maintain robustness in the presence of noisy reward feedback. We contend that alignment techniques such as \textbf{WEAK-TO-STRONG} \citep{burns2023weak} will substantially transform existing learning algorithm paradigms.

\paragraph{Reward Function} The reward function serves as the principal training signal in reinforcement learning. Typically, this role is fulfilled by a neural reward model. We identify three key areas for the advancement of reward models: 1) \textbf{Enhancing the generalization capabilities of the reward model.} It is critical that the reward model generalizes robustly to both out-of-distribution prompts and more sophisticated decoded outputs; failure to do so could result in reinforcement learning merely reinforcing the initial LLM output distributions rather than enabling significant capability improvements; 2) \textbf{Accounting for the uncertainty inherent in the reward model.} Modeling this uncertainty may facilitate the integration between weaker reward models and learning algorithms that transition from weak to strong performance; 3) \textbf{Constructing high-quality process reward models efficiently} so as to deliver precise, granular training signals throughout the reasoning trajectory \citep{lightman2023let,wang2023math}.

\section{Conclusion, Limitation, and Future Work}

In this work, we introduce DeepSeekMath, a model that surpasses all available open-source systems on the MATH competition benchmark and comes close to matching the achievements of proprietary alternatives. DeepSeekMath~is built upon DeepSeek-Coder-v1.5 7B, which is subjected to continual training over 500B tokens; a substantial fraction of this training set—120B tokens—consists of mathematical content derived from Common Crawl. Through a comprehensive ablation analysis, we demonstrate that web page data is a particularly rich source for quality mathematics, while arXiv data does not offer the expected gains. We propose Group Relative Policy Optimization (GRPO), which adapts Proximal Policy Optimization (PPO) to deliver enhanced mathematical reasoning performance with reduced memory overhead. Experimental evaluation indicates that GRPO confers notable improvements even when DeepSeekMath-Instruct 7B already achieves strong results on standard benchmarks. Finally, we introduce a unified framework for understanding various related methodologies and outline multiple promising avenues for further enhancement of reinforcement learning effectiveness.

While DeepSeekMath~performs strongly on benchmarks involving quantitative reasoning, its abilities in geometry and theorem-proving remain less robust compared to proprietary models. For example, during preliminary testing, the model was unable to solve questions involving triangles and ellipses, suggesting potential data selection biases during pre-training and fine-tuning stages. Moreover, due to constraints in model size, DeepSeekMath~displays inferior few-shot learning capabilities relative to GPT-4. Unlike GPT-4, which benefits from few-shot prompts, DeepSeekMath~demonstrates little variation between zero-shot and few-shot results. Going forward, we plan to enhance our data selection pipeline to produce a higher-quality pre-training dataset. Furthermore, we will investigate novel approaches for more effective reinforcement learning of LLMs, as discussed in Section \ref{sec:effec_rl}.

\end{document}